\section{The physical cylinder exponent}\label{c7:cylinder-exponent}

The operator estimate in Section~\ref{c7:compact} controls the normalized
quantity $R_m$ from \eqref{j:normalized-cylinder}. We now insert the
homogeneous scalar normalization and return to the positive cylinder
partition. Here $A,B$ are the two intervals on the balanced horizontal
cut of the introduction, so $G_{A,B}=G_m$. Each has size $m$, and $m$
tends to infinity through even integers.

For $X=(1,\ldots,1)$ of size $2m$, \eqref{j:homogeneous-normalization} gives
\[
 N(X)=(2p_0/3)^{m(2m-1)}(2/3)^m m^{5/48+o(1)},
 \qquad p_0=2+\sqrt2,
\]
where the absolute value would suffice. Together with $P_0$ and
$\kappa_\infty$, this gives
\begin{equation}\label{j:physical-scalar-normalization}
 e^{-m^2\kappa_\infty}\frac{N(X)^2}{P_0}=m^{5/24+o(1)}.
\end{equation}
Put $\Theta=-2\theta$, $g=2\cos\Theta$, and
\[
 p_*(\Theta)=\frac16-\frac{2\Theta^2}{3\pi^2}.
\]
In \eqref{j:balanced-trace-estimate}, $u=u_0$ and
$L=\log m/(\kappa_1b)+O(1)$, with $\kappa_1=\pi/3$. The contribution of
the interval follows from the exact exponential expression
\[
 \begin{gathered}
 \frac{u}{r_E}=\exp\!\left(\frac{b\pi}{16}
                    -\frac{4b\theta^2}{9\pi}\right),\\
 \frac{2}{\kappa_1b}\left(\frac{b\pi}{16}
                    -\frac{4b\theta^2}{9\pi}\right)
 =\frac38-\frac{8\theta^2}{3\pi^2}.
 \end{gathered}
\]
Subtracting the exponent $5/24$ in
\eqref{j:physical-scalar-normalization} proves the locally uniform bound
\begin{equation}\label{j:cylinder-complex-bound}
 |G_{A,B}(2\cos\Theta)|
 \le m^{\Re p_*(\Theta)+o(1)},\qquad |\Re\Theta|<\pi.
\end{equation}

The exact normalization at $g=0$ supplies the logarithmic derivative
bound. At $\theta_i=-\pi/4$, equivalently $\Theta=\pi/2$, the partition
is $G_{A,B}(0)=1$ and $p_*(\pi/2)=0$. Thus it saturates the exponent in
\eqref{j:cylinder-complex-bound}. In particular the analytic coefficient
$\mathcal F_\zeta(a,\theta_i)$ in
\eqref{j:balanced-trace-profile} cannot vanish at any point of
$[1,e^{\kappa_1b}]$. Indeed even integers $m$ can approximate any point
of this interval by $a_m$ with error exponentially small in $L$, from
inside at the endpoints. A zero, together with the exponentially small
error in \eqref{j:balanced-trace-profile}, would make
$|R_m(u/r_E)^{-2L}|$ exponentially small in $L$. This contradicts
\eqref{j:physical-scalar-normalization} and $G_{A,B}(0)=1$, which give
$|R_m(u/r_E)^{-2L}|=m^{o(1)}$ there.

Compactness of the $a$ interval now gives a parameter neighborhood of
$\theta_i$ on which the analytic factor stays nonzero and its fixed-order
logarithmic derivatives are bounded. The analytic error in
\eqref{j:balanced-trace-profile} has the same local derivative control
on a smaller neighborhood by Cauchy estimates.
Differentiating the resulting expression for $G_{A,B}$ along the local
branch $\Theta(g)$ near $g=0$ therefore yields
\begin{equation}\label{j:cylinder-pair-bound}
 -\partial_g^2\log G_{A,B}(0)=O(\log m).
\end{equation}

It remains to obtain the matching lower bound at $g=2$. We claim
\begin{equation}\label{j:cylinder-positive-exponent}
 G_{A,B}(2)=m^{1/6+o(1)}.
\end{equation}
The upper bound is \eqref{j:cylinder-complex-bound}. If the lower bound
failed along a sequence by a fixed exponent gap, positivity of the
polynomial coefficients would give $|G_{A,B}(g)|\le G_{A,B}(2)$ for
$|g|\le2$. Choose a sufficiently small nonzero real $\Theta_0$ and then
a disk about it on which $|2\cos\Theta|<2$ and $\Re p_*(\Theta)$ is
sufficiently close to $1/6$. On this disk the analytic function
\[
 G_{A,B}(2\cos\Theta)m^{-p_*(\Theta)}
\]
would have modulus at most $m^{-\delta}$ for some fixed $\delta>0$. Its local upper
bound elsewhere in the strip is subpower by
\eqref{j:cylinder-complex-bound}. The three-circle inequality on finitely
many overlapping disks along the interval from $\Theta_0$ to $\pi/2$
propagates a positive fraction of the exponent gap to $\pi/2$. This
contradicts the exact value one there. This proves
\eqref{j:cylinder-positive-exponent}, and with
\eqref{j:cylinder-complex-bound}--\eqref{j:cylinder-pair-bound} proves
the cylinder assertions of Theorem~\ref{j:cylinder}.
